\documentclass[12pt,a4paper]{article}

\usepackage[utf8]{inputenc}
\usepackage[T2A]{fontenc}
\usepackage[russian]{babel}
\usepackage{amsmath,amssymb,amsthm}
\usepackage{listings}
\usepackage{xcolor}
\usepackage{geometry}
\usepackage{array}
\usepackage{booktabs}
\usepackage{enumitem}
\usepackage{hyperref}

\geometry{left=2cm,right=2cm,top=2cm,bottom=2cm}

\lstset{
    basicstyle=\ttfamily\small,
    keywordstyle=\color{blue},
    commentstyle=\color{gray},
    stringstyle=\color{red!60!black},
    numbers=left,
    numberstyle=\tiny,
    stepnumber=1,
    numbersep=8pt,
    frame=single,
    breaklines=true,
    tabsize=4,
    showstringspaces=false,
    literate=
      {А}{{\CYRA}}1
      {Б}{{\CYRB}}1
      {В}{{\CYRV}}1
      {Г}{{\CYRG}}1
      {Д}{{\CYRD}}1
      {Е}{{\CYRE}}1
      {Ё}{{\CYRYO}}1
      {Ж}{{\CYRZH}}1
      {З}{{\CYRZ}}1
      {И}{{\CYRI}}1
      {Й}{{\CYRISHRT}}1
      {К}{{\CYRK}}1
      {Л}{{\CYRL}}1
      {М}{{\CYRM}}1
      {Н}{{\CYRN}}1
      {О}{{\CYRO}}1
      {П}{{\CYRP}}1
      {Р}{{\CYRR}}1
      {С}{{\CYRS}}1
      {Т}{{\CYRT}}1
      {У}{{\CYRU}}1
      {Ф}{{\CYRF}}1
      {Х}{{\CYRH}}1
      {Ц}{{\CYRC}}1
      {Ч}{{\CYRCH}}1
      {Ш}{{\CYRSH}}1
      {Щ}{{\CYRSHCH}}1
      {Ъ}{{\CYRHRDSN}}1
      {Ы}{{\CYRERY}}1
      {Ь}{{\CYRSFTSN}}1
      {Э}{{\CYREREV}}1
      {Ю}{{\CYRYU}}1
      {Я}{{\CYRYA}}1
      {а}{{\cyra}}1
      {б}{{\cyrb}}1
      {в}{{\cyrv}}1
      {г}{{\cyrg}}1
      {д}{{\cyrd}}1
      {е}{{\cyre}}1
      {ё}{{\cyryo}}1
      {ж}{{\cyrzh}}1
      {з}{{\cyrz}}1
      {и}{{\cyri}}1
      {й}{{\cyrishrt}}1
      {к}{{\cyrk}}1
      {л}{{\cyrl}}1
      {м}{{\cyrm}}1
      {н}{{\cyrn}}1
      {о}{{\cyro}}1
      {п}{{\cyrp}}1
      {р}{{\cyrr}}1
      {с}{{\cyrs}}1
      {т}{{\cyrt}}1
      {у}{{\cyru}}1
      {ф}{{\cyrf}}1
      {х}{{\cyrh}}1
      {ц}{{\cyrc}}1
      {ч}{{\cyrch}}1
      {ш}{{\cyrsh}}1
      {щ}{{\cyrshch}}1
      {ъ}{{\cyrhrdsn}}1
      {ы}{{\cyrery}}1
      {ь}{{\cyrsftsn}}1
      {э}{{\cyrerev}}1
      {ю}{{\cyryu}}1
      {я}{{\cyrya}}1
}

\title{Билет №67 \\ \small Операционные средства управления процессами при их реализации на параллельных и распределенных вычислительных системах и сетях: стандарты и программные средства PVM, MPI, OpenMP, POSIX. (ИТМО)}
\author{Сериков Владислав Александрович}
\date{16 мая 2026}

\begin{document}

\maketitle
\tableofcontents
\newpage

\section{Введение}

Параллельные и распределённые вычислительные системы используются для ускорения вычислений, повышения отказоустойчивости, обработки больших объёмов данных и эффективного использования совокупных ресурсов множества процессоров, узлов и сетей.

\textbf{Параллельная вычислительная система} --- система, в которой несколько вычислительных устройств совместно решают одну задачу. Обычно предполагается тесная связь между вычислителями и сравнительно малая задержка обмена.

\textbf{Распределённая вычислительная система} --- совокупность автономных вычислительных узлов, связанных сетью и взаимодействующих посредством обмена сообщениями. Узлы могут быть географически удалены, иметь разные архитектуры и операционные системы.

\textbf{Процесс} --- экземпляр выполняющейся программы, обладающий собственным адресным пространством, набором ресурсов и состоянием выполнения.

\textbf{Поток} --- единица выполнения внутри процесса. Потоки одного процесса разделяют адресное пространство, но имеют собственные регистры, стек и состояние планирования.

\textbf{Управление процессами} в параллельных и распределённых системах включает:
\begin{itemize}
    \item создание, запуск и завершение процессов;
    \item распределение процессов по узлам;
    \item синхронизацию процессов и потоков;
    \item обмен сообщениями;
    \item управление ресурсами;
    \item обработку ошибок и отказов;
    \item балансировку нагрузки;
    \item поддержку переносимости программ.
\end{itemize}

К основным программным средствам и стандартам, применяемым для этих целей, относятся:
\begin{itemize}
    \item PVM --- Parallel Virtual Machine;
    \item MPI --- Message Passing Interface;
    \item OpenMP --- Open Multi-Processing;
    \item POSIX --- Portable Operating System Interface.
\end{itemize}

\section{Архитектуры параллельных и распределённых систем}

\subsection{Классификация по памяти}

\subsubsection{Системы с общей памятью}

В системах с общей памятью все процессоры имеют доступ к единому адресному пространству.

\[
P_1, P_2, \ldots, P_n \longrightarrow M
\]

где $P_i$ --- процессоры, $M$ --- общая память.

Преимущества:
\begin{itemize}
    \item простая модель программирования;
    \item быстрый обмен данными через общие переменные;
    \item удобная реализация потоков.
\end{itemize}

Недостатки:
\begin{itemize}
    \item ограниченная масштабируемость;
    \item необходимость синхронизации доступа к памяти;
    \item возможные состояния гонки.
\end{itemize}

Типичные средства программирования:
\begin{itemize}
    \item POSIX Threads;
    \item OpenMP;
    \item C++ Threads.
\end{itemize}

\subsubsection{Системы с распределённой памятью}

В системах с распределённой памятью каждый процессор или узел имеет собственную локальную память. Обмен данными выполняется через сеть.

\[
(P_1, M_1) \leftrightarrow (P_2, M_2) \leftrightarrow \ldots \leftrightarrow (P_n, M_n)
\]

Преимущества:
\begin{itemize}
    \item высокая масштабируемость;
    \item возможность использования кластеров;
    \item отсутствие общего адресного пространства снижает количество ошибок, связанных с гонками за память.
\end{itemize}

Недостатки:
\begin{itemize}
    \item необходимость явного обмена сообщениями;
    \item зависимость производительности от сети;
    \item более сложное программирование.
\end{itemize}

Типичные средства программирования:
\begin{itemize}
    \item MPI;
    \item PVM.
\end{itemize}

\subsubsection{Гибридные системы}

Современные суперкомпьютеры часто имеют гибридную архитектуру: несколько узлов связаны сетью, а внутри каждого узла есть многоядерный процессор с общей памятью.

Типичная модель программирования:

\[
\text{MPI между узлами} + \text{OpenMP внутри узла}.
\]

\section{Понятие процесса в параллельных вычислениях}

\subsection{Состояния процесса}

В операционной системе процесс обычно находится в одном из следующих состояний:
\begin{itemize}
    \item \textbf{создан} --- процесс создан, но ещё не запущен;
    \item \textbf{готов} --- процесс ожидает выделения процессора;
    \item \textbf{выполняется} --- процесс исполняется на процессоре;
    \item \textbf{ожидает} --- процесс ожидает события, например завершения ввода-вывода или получения сообщения;
    \item \textbf{завершён} --- процесс закончил выполнение.
\end{itemize}

Схематически:

\[
\text{создан} \rightarrow \text{готов} \rightarrow \text{выполняется} \rightarrow \text{завершён}.
\]

Также возможен переход:

\[
\text{выполняется} \rightarrow \text{ожидает} \rightarrow \text{готов}.
\]

\subsection{Процессы и потоки}

\begin{center}
\begin{tabular}{|p{5cm}|p{5cm}|p{5cm}|}
\hline
\textbf{Свойство} & \textbf{Процесс} & \textbf{Поток} \\
\hline
Адресное пространство & Собственное & Общее с другими потоками процесса \\
\hline
Стоимость создания & Относительно высокая & Относительно низкая \\
\hline
Обмен данными & Через IPC, файлы, сокеты, сообщения & Через общие переменные \\
\hline
Изоляция & Высокая & Низкая \\
\hline
Типичные средства & MPI, PVM, POSIX process API & POSIX Threads, OpenMP \\
\hline
\end{tabular}
\end{center}

\section{Основные модели параллельного программирования}

\subsection{Модель передачи сообщений}

В модели передачи сообщений каждый процесс имеет собственную память. Для обмена данными процессы явно отправляют и принимают сообщения.

Основные операции:
\begin{itemize}
    \item \texttt{send} --- отправка сообщения;
    \item \texttt{receive} --- получение сообщения;
    \item \texttt{broadcast} --- широковещательная рассылка;
    \item \texttt{reduce} --- коллективное вычисление;
    \item \texttt{barrier} --- барьерная синхронизация.
\end{itemize}

Данная модель используется в MPI и PVM.

\subsection{Модель общей памяти}

В модели общей памяти потоки взаимодействуют через общие переменные. Для предотвращения ошибок применяются:
\begin{itemize}
    \item мьютексы;
    \item семафоры;
    \item критические секции;
    \item атомарные операции;
    \item барьеры.
\end{itemize}

Данная модель используется в OpenMP и POSIX Threads.

\subsection{SPMD-модель}

\textbf{SPMD} --- Single Program Multiple Data.

Все процессы выполняют одну и ту же программу, но обрабатывают разные части данных.

Пример:
\[
\text{процесс } i \text{ обрабатывает блок массива } A_i.
\]

SPMD широко применяется в MPI-программах.

\section{PVM}

\subsection{Общая характеристика}

\textbf{PVM} --- Parallel Virtual Machine --- программная система, позволяющая объединять набор разнородных компьютеров в единую виртуальную параллельную машину.

PVM ориентирована на распределённые системы, состоящие из узлов с разными архитектурами, операционными системами и характеристиками сети.

Главная идея PVM состоит в том, что физически разнородная сеть компьютеров представляется программисту как один параллельный вычислительный ресурс.

\subsection{Архитектура PVM}

PVM состоит из:
\begin{itemize}
    \item демона \texttt{pvmd}, работающего на каждом узле;
    \item библиотеки PVM, подключаемой к пользовательским программам;
    \item консольных утилит управления виртуальной машиной;
    \item механизма обмена сообщениями между задачами.
\end{itemize}

Каждая задача PVM имеет уникальный идентификатор --- \textbf{TID}.

Схема работы:

\[
\text{программа} \leftrightarrow \text{PVM-библиотека} \leftrightarrow \text{pvmd} \leftrightarrow \text{сеть} \leftrightarrow \text{pvmd}.
\]

\subsection{Основные функции PVM}

\begin{center}
\begin{tabular}{|p{5cm}|p{9cm}|}
\hline
\textbf{Функция} & \textbf{Назначение} \\
\hline
\texttt{pvm\_spawn} & Запуск новых задач \\
\hline
\texttt{pvm\_mytid} & Получение идентификатора текущей задачи \\
\hline
\texttt{pvm\_parent} & Получение идентификатора родительской задачи \\
\hline
\texttt{pvm\_initsend} & Инициализация буфера отправки \\
\hline
\texttt{pvm\_pkint}, \texttt{pvm\_pkdouble} & Упаковка данных в сообщение \\
\hline
\texttt{pvm\_send} & Отправка сообщения \\
\hline
\texttt{pvm\_recv} & Получение сообщения \\
\hline
\texttt{pvm\_upkint}, \texttt{pvm\_upkdouble} & Распаковка данных \\
\hline
\texttt{pvm\_exit} & Завершение участия задачи в PVM \\
\hline
\end{tabular}
\end{center}

\subsection{Алгоритм запуска задач в PVM}

\textbf{Задача:} запустить несколько рабочих процессов из главного процесса.

\textbf{Шаги алгоритма:}
\begin{enumerate}
    \item Главный процесс подключается к виртуальной машине PVM.
    \item Получает собственный идентификатор \texttt{TID}.
    \item Вызывает \texttt{pvm\_spawn} для запуска рабочих задач.
    \item Передаёт рабочим задачам данные через сообщения.
    \item Рабочие задачи выполняют вычисления.
    \item Рабочие задачи отправляют результаты главному процессу.
    \item Главный процесс собирает результаты.
    \item Все процессы завершают работу с PVM.
\end{enumerate}

Иллюстрация:

\[
\text{Master}
\rightarrow
\begin{cases}
\text{Worker}_1,\\
\text{Worker}_2,\\
\ldots\\
\text{Worker}_n
\end{cases}
\rightarrow
\text{Master}.
\]

\subsection{Минимальный пример PVM}

\subsubsection{Главная программа}

\begin{lstlisting}[language=C]
#include <stdio.h>
#include "pvm3.h"

int main() {
    int tid;
    int child;
    int value = 42;
    int result;

    tid = pvm_mytid();

    pvm_spawn("worker", NULL, PvmTaskDefault, "", 1, &child);

    pvm_initsend(PvmDataDefault);
    pvm_pkint(&value, 1, 1);
    pvm_send(child, 1);

    pvm_recv(child, 2);
    pvm_upkint(&result, 1, 1);

    printf("Result = %d\n", result);

    pvm_exit();
    return 0;
}
\end{lstlisting}

\subsubsection{Рабочая программа}

\begin{lstlisting}[language=C]
#include "pvm3.h"

int main() {
    int parent;
    int value;

    parent = pvm_parent();

    pvm_recv(parent, 1);
    pvm_upkint(&value, 1, 1);

    value = value * 2;

    pvm_initsend(PvmDataDefault);
    pvm_pkint(&value, 1, 1);
    pvm_send(parent, 2);

    pvm_exit();
    return 0;
}
\end{lstlisting}

\subsection{Достоинства и недостатки PVM}

\textbf{Достоинства:}
\begin{itemize}
    \item поддержка гетерогенных систем;
    \item возможность динамического добавления и удаления узлов;
    \item удобная модель виртуальной машины;
    \item относительно простая организация распределённых вычислений.
\end{itemize}

\textbf{Недостатки:}
\begin{itemize}
    \item меньшая стандартизованность по сравнению с MPI;
    \item в современных HPC-системах используется значительно реже MPI;
    \item более низкая производительность в ряде реализаций;
    \item ограниченная поддержка современных суперкомпьютерных архитектур.
\end{itemize}

\section{MPI}

\subsection{Общая характеристика}

\textbf{MPI} --- Message Passing Interface --- стандарт интерфейса передачи сообщений для параллельного программирования в системах с распределённой памятью.

MPI является не конкретной библиотекой, а спецификацией. Существуют разные реализации MPI, например:
\begin{itemize}
    \item MPICH;
    \item Open MPI;
    \item Intel MPI;
    \item MVAPICH.
\end{itemize}

MPI используется в кластерах, суперкомпьютерах и гибридных системах.

\subsection{Основные понятия MPI}

\subsubsection{Коммуникатор}

\textbf{Коммуникатор} --- объект, определяющий группу процессов, которые могут взаимодействовать друг с другом.

Стандартный коммуникатор:
\[
\texttt{MPI\_COMM\_WORLD}
\]
содержит все процессы, запущенные в рамках MPI-программы.

\subsubsection{Ранг процесса}

\textbf{Ранг} --- целочисленный номер процесса внутри коммуникатора.

Если в коммуникаторе $p$ процессов, то ранги имеют значения:
\[
0,1,\ldots,p-1.
\]

\subsubsection{Сообщение}

Сообщение в MPI задаётся:
\begin{itemize}
    \item адресом буфера;
    \item количеством элементов;
    \item типом данных;
    \item рангом отправителя или получателя;
    \item тегом;
    \item коммуникатором.
\end{itemize}

\subsection{Базовая структура MPI-программы}

\begin{lstlisting}[language=C]
#include <mpi.h>
#include <stdio.h>

int main(int argc, char **argv) {
    int rank, size;

    MPI_Init(&argc, &argv);

    MPI_Comm_rank(MPI_COMM_WORLD, &rank);
    MPI_Comm_size(MPI_COMM_WORLD, &size);

    printf("Hello from process %d of %d\n", rank, size);

    MPI_Finalize();
    return 0;
}
\end{lstlisting}

Компиляция и запуск, например:

\begin{lstlisting}
mpicc hello.c -o hello
mpirun -np 4 ./hello
\end{lstlisting}

\subsection{Точечные обмены}

\subsubsection{Блокирующая отправка и получение}

Основные функции:
\begin{itemize}
    \item \texttt{MPI\_Send};
    \item \texttt{MPI\_Recv}.
\end{itemize}

Пример:

\begin{lstlisting}[language=C]
if (rank == 0) {
    int x = 100;
    MPI_Send(&x, 1, MPI_INT, 1, 0, MPI_COMM_WORLD);
} else if (rank == 1) {
    int x;
    MPI_Recv(&x, 1, MPI_INT, 0, 0, MPI_COMM_WORLD,
             MPI_STATUS_IGNORE);
    printf("Received %d\n", x);
}
\end{lstlisting}

\subsubsection{Неблокирующие обмены}

Основные функции:
\begin{itemize}
    \item \texttt{MPI\_Isend};
    \item \texttt{MPI\_Irecv};
    \item \texttt{MPI\_Wait};
    \item \texttt{MPI\_Test}.
\end{itemize}

Неблокирующие операции позволяют перекрывать вычисления и коммуникации.

\begin{lstlisting}[language=C]
MPI_Request req;
MPI_Isend(&x, 1, MPI_INT, dest, tag, MPI_COMM_WORLD, &req);

/* Здесь можно выполнять вычисления */

MPI_Wait(&req, MPI_STATUS_IGNORE);
\end{lstlisting}

\subsection{Коллективные операции MPI}

Коллективные операции выполняются всеми процессами коммуникатора.

\begin{center}
\begin{tabular}{|p{5cm}|p{9cm}|}
\hline
\textbf{Функция} & \textbf{Назначение} \\
\hline
\texttt{MPI\_Bcast} & Рассылка данных от одного процесса всем остальным \\
\hline
\texttt{MPI\_Scatter} & Разбиение массива и рассылка частей разным процессам \\
\hline
\texttt{MPI\_Gather} & Сбор частей массива на одном процессе \\
\hline
\texttt{MPI\_Reduce} & Выполнение редукции, например суммы или максимума \\
\hline
\texttt{MPI\_Allreduce} & Редукция с выдачей результата всем процессам \\
\hline
\texttt{MPI\_Barrier} & Барьерная синхронизация \\
\hline
\end{tabular}
\end{center}

\subsection{Алгоритм параллельного суммирования массива в MPI}

\textbf{Задача:} вычислить сумму элементов массива $A$ длины $N$ на $p$ процессах.

\textbf{Идея:} массив делится на $p$ блоков, каждый процесс суммирует свой блок, затем частичные суммы складываются.

\textbf{Шаги алгоритма:}
\begin{enumerate}
    \item Процесс с рангом $0$ создаёт исходный массив $A$.
    \item Массив разбивается на $p$ частей.
    \item С помощью \texttt{MPI\_Scatter} каждая часть передаётся соответствующему процессу.
    \item Каждый процесс вычисляет локальную сумму:
    \[
    s_i = \sum_{j \in B_i} A_j,
    \]
    где $B_i$ --- блок элементов, назначенный процессу $i$.
    \item С помощью \texttt{MPI\_Reduce} локальные суммы объединяются:
    \[
    S = \sum_{i=0}^{p-1} s_i.
    \]
    \item Процесс $0$ получает итоговую сумму $S$.
\end{enumerate}

Иллюстрация для $p=4$:

\[
A = [A_0|A_1|A_2|A_3]
\]

\[
P_0 \leftarrow A_0,\quad
P_1 \leftarrow A_1,\quad
P_2 \leftarrow A_2,\quad
P_3 \leftarrow A_3.
\]

\[
S = s_0+s_1+s_2+s_3.
\]

\subsection{Минимальный пример MPI-суммирования}

\begin{lstlisting}[language=C]
#include <mpi.h>
#include <stdio.h>

int main(int argc, char **argv) {
    int rank, size;
    int local_value;
    int local_sum;
    int global_sum;

    MPI_Init(&argc, &argv);

    MPI_Comm_rank(MPI_COMM_WORLD, &rank);
    MPI_Comm_size(MPI_COMM_WORLD, &size);

    local_value = rank + 1;
    local_sum = local_value;

    MPI_Reduce(&local_sum, &global_sum, 1, MPI_INT,
               MPI_SUM, 0, MPI_COMM_WORLD);

    if (rank == 0) {
        printf("Sum = %d\n", global_sum);
    }

    MPI_Finalize();
    return 0;
}
\end{lstlisting}

Если запустить программу на четырёх процессах, значения будут:
\[
1,2,3,4,
\]
а сумма:
\[
1+2+3+4=10.
\]

\subsection{Синхронизация в MPI}

Основное средство глобальной синхронизации --- барьер:

\begin{lstlisting}[language=C]
MPI_Barrier(MPI_COMM_WORLD);
\end{lstlisting}

Барьер означает, что ни один процесс не продолжит выполнение после барьера, пока все процессы коммуникатора не достигнут этой точки.

\subsection{Достоинства и недостатки MPI}

\textbf{Достоинства:}
\begin{itemize}
    \item высокий уровень стандартизации;
    \item переносимость между суперкомпьютерами и кластерами;
    \item высокая производительность;
    \item поддержка коллективных операций;
    \item поддержка сложных топологий и пользовательских типов данных;
    \item пригодность для масштабирования на тысячи и миллионы процессов.
\end{itemize}

\textbf{Недостатки:}
\begin{itemize}
    \item необходимость явного управления обменами;
    \item сложность отладки;
    \item риск взаимных блокировок;
    \item усложнение программ при нерегулярных обменах.
\end{itemize}

\section{OpenMP}

\subsection{Общая характеристика}

\textbf{OpenMP} --- стандарт API для параллельного программирования в системах с общей памятью. Он основан на использовании директив компилятора, функций библиотеки и переменных окружения.

OpenMP поддерживает языки:
\begin{itemize}
    \item C;
    \item C++;
    \item Fortran.
\end{itemize}

Главная идея OpenMP --- постепенное распараллеливание последовательной программы с помощью директив.

\subsection{Модель выполнения OpenMP}

OpenMP использует модель \textbf{fork-join}.

\begin{enumerate}
    \item Программа начинается как один главный поток.
    \item При входе в параллельную область главный поток создаёт группу потоков.
    \item Потоки совместно выполняют код параллельной области.
    \item В конце параллельной области потоки синхронизируются.
    \item Выполнение снова продолжается одним главным потоком.
\end{enumerate}

Иллюстрация:

\[
\text{Master}
\rightarrow
\begin{cases}
T_0\\
T_1\\
T_2\\
T_3
\end{cases}
\rightarrow
\text{Master}.
\]

\subsection{Минимальный пример OpenMP}

\begin{lstlisting}[language=C]
#include <stdio.h>
#include <omp.h>

int main() {
#pragma omp parallel
    {
        int id = omp_get_thread_num();
        int n = omp_get_num_threads();
        printf("Hello from thread %d of %d\n", id, n);
    }

    return 0;
}
\end{lstlisting}

Компиляция:

\begin{lstlisting}
gcc -fopenmp hello.c -o hello
\end{lstlisting}

\subsection{Основные директивы OpenMP}

\begin{center}
\begin{tabular}{|p{5cm}|p{9cm}|}
\hline
\textbf{Директива} & \textbf{Назначение} \\
\hline
\texttt{\#pragma omp parallel} & Создание параллельной области \\
\hline
\texttt{\#pragma omp for} & Распределение итераций цикла между потоками \\
\hline
\texttt{\#pragma omp parallel for} & Совмещённое создание потоков и распараллеливание цикла \\
\hline
\texttt{\#pragma omp sections} & Разделение независимых секций кода между потоками \\
\hline
\texttt{\#pragma omp single} & Выполнение блока одним потоком \\
\hline
\texttt{\#pragma omp critical} & Критическая секция \\
\hline
\texttt{\#pragma omp atomic} & Атомарная операция \\
\hline
\texttt{\#pragma omp barrier} & Барьерная синхронизация \\
\hline
\texttt{\#pragma omp task} & Создание задачи \\
\hline
\end{tabular}
\end{center}

\subsection{Параллельный цикл}

\begin{lstlisting}[language=C]
#pragma omp parallel for
for (int i = 0; i < n; i++) {
    c[i] = a[i] + b[i];
}
\end{lstlisting}

Итерации цикла распределяются между потоками.

\subsection{Алгоритм параллельного суммирования в OpenMP}

\textbf{Задача:} вычислить сумму элементов массива $A$.

\textbf{Шаги алгоритма:}
\begin{enumerate}
    \item Создать переменную \texttt{sum}.
    \item Запустить параллельный цикл.
    \item Разделить итерации между потоками.
    \item Каждый поток вычисляет частичную сумму.
    \item OpenMP объединяет частичные суммы с помощью редукции.
    \item Получается итоговая сумма.
\end{enumerate}

Математически:
\[
S = \sum_{i=0}^{N-1} A_i.
\]

Если потоки $T_0,\ldots,T_{p-1}$ получают блоки $B_0,\ldots,B_{p-1}$, то:
\[
S_k = \sum_{i \in B_k} A_i,
\]
\[
S = \sum_{k=0}^{p-1} S_k.
\]

\subsection{Минимальный пример редукции OpenMP}

\begin{lstlisting}[language=C]
#include <stdio.h>
#include <omp.h>

int main() {
    int n = 100;
    int sum = 0;

#pragma omp parallel for reduction(+:sum)
    for (int i = 1; i <= n; i++) {
        sum += i;
    }

    printf("Sum = %d\n", sum);
    return 0;
}
\end{lstlisting}

Результат:
\[
1+2+\ldots+100 = \frac{100 \cdot 101}{2}=5050.
\]

\subsection{Переменные в OpenMP}

OpenMP различает:
\begin{itemize}
    \item \textbf{shared} --- общие переменные;
    \item \textbf{private} --- локальные переменные потока;
    \item \textbf{firstprivate} --- локальная копия с начальным значением;
    \item \textbf{lastprivate} --- сохранение последнего значения после параллельной области;
    \item \textbf{reduction} --- переменная редукции.
\end{itemize}

Пример:

\begin{lstlisting}[language=C]
#pragma omp parallel for private(i) shared(a, b, c)
for (i = 0; i < n; i++) {
    c[i] = a[i] + b[i];
}
\end{lstlisting}

\subsection{Планирование циклов}

Директива \texttt{schedule} задаёт способ распределения итераций.

\begin{center}
\begin{tabular}{|p{5cm}|p{9cm}|}
\hline
\textbf{Вид планирования} & \textbf{Описание} \\
\hline
\texttt{static} & Итерации распределяются заранее равными блоками \\
\hline
\texttt{dynamic} & Потоки получают новые блоки по мере освобождения \\
\hline
\texttt{guided} & Размер блоков постепенно уменьшается \\
\hline
\texttt{runtime} & Способ задаётся переменной окружения \texttt{OMP\_SCHEDULE} \\
\hline
\end{tabular}
\end{center}

Пример:

\begin{lstlisting}[language=C]
#pragma omp parallel for schedule(dynamic, 10)
for (int i = 0; i < n; i++) {
    work(i);
}
\end{lstlisting}

\subsection{Состояния гонки}

\textbf{Состояние гонки} возникает, когда несколько потоков одновременно обращаются к общей переменной, хотя бы один поток выполняет запись, а результат зависит от порядка выполнения.

Пример ошибки:

\begin{lstlisting}[language=C]
int counter = 0;

#pragma omp parallel for
for (int i = 0; i < 1000; i++) {
    counter++;
}
\end{lstlisting}

Операция \texttt{counter++} не является атомарной.

Исправление через \texttt{atomic}:

\begin{lstlisting}[language=C]
#pragma omp parallel for
for (int i = 0; i < 1000; i++) {
#pragma omp atomic
    counter++;
}
\end{lstlisting}

Но для суммирования лучше использовать \texttt{reduction}.

\subsection{Достоинства и недостатки OpenMP}

\textbf{Достоинства:}
\begin{itemize}
    \item простое добавление параллелизма в существующий код;
    \item удобство для многоядерных систем с общей памятью;
    \item поддержка директив, библиотечных функций и переменных окружения;
    \item сравнительно низкий порог входа.
\end{itemize}

\textbf{Недостатки:}
\begin{itemize}
    \item ограниченность системами с общей памятью;
    \item риск состояний гонки;
    \item сложность контроля производительности на NUMA-системах;
    \item зависимость от компилятора и его поддержки стандарта.
\end{itemize}

\section{POSIX}

\subsection{Общая характеристика}

\textbf{POSIX} --- Portable Operating System Interface --- семейство стандартов, определяющих интерфейсы операционной системы, прежде всего для UNIX-подобных систем.

POSIX описывает:
\begin{itemize}
    \item API системных вызовов;
    \item работу с процессами;
    \item работу с потоками;
    \item файловую систему;
    \item сигналы;
    \item межпроцессное взаимодействие;
    \item средства синхронизации;
    \item командную оболочку и утилиты.
\end{itemize}

POSIX важен для параллельных и распределённых систем, так как обеспечивает переносимые механизмы управления процессами и потоками.

\subsection{Управление процессами в POSIX}

Основные функции:
\begin{center}
\begin{tabular}{|p{5cm}|p{9cm}|}
\hline
\textbf{Функция} & \textbf{Назначение} \\
\hline
\texttt{fork()} & Создание нового процесса \\
\hline
\texttt{exec()} & Замена образа процесса новой программой \\
\hline
\texttt{wait()} / \texttt{waitpid()} & Ожидание завершения дочернего процесса \\
\hline
\texttt{exit()} & Завершение процесса \\
\hline
\texttt{getpid()} & Получение PID текущего процесса \\
\hline
\texttt{getppid()} & Получение PID родительского процесса \\
\hline
\end{tabular}
\end{center}

\subsection{Алгоритм создания процесса в POSIX}

\textbf{Шаги алгоритма:}
\begin{enumerate}
    \item Родительский процесс вызывает \texttt{fork()}.
    \item Операционная система создаёт копию процесса.
    \item В родительском процессе \texttt{fork()} возвращает PID дочернего процесса.
    \item В дочернем процессе \texttt{fork()} возвращает $0$.
    \item Дочерний процесс может выполнить другую программу через \texttt{exec()}.
    \item Родительский процесс ожидает завершения дочернего через \texttt{wait()}.
\end{enumerate}

Иллюстрация:

\[
\text{Parent}
\xrightarrow{\texttt{fork()}}
\begin{cases}
\text{Parent},\\
\text{Child}.
\end{cases}
\]

\subsection{Минимальный пример \texttt{fork}}

\begin{lstlisting}[language=C]
#include <stdio.h>
#include <unistd.h>
#include <sys/wait.h>

int main() {
    pid_t pid = fork();

    if (pid == 0) {
        printf("Child process, pid = %d\n", getpid());
    } else {
        wait(NULL);
        printf("Parent process, child pid = %d\n", pid);
    }

    return 0;
}
\end{lstlisting}

\subsection{POSIX Threads}

\textbf{POSIX Threads}, или \textbf{pthreads}, --- стандартный POSIX-интерфейс для создания и управления потоками.

Основные функции:
\begin{center}
\begin{tabular}{|p{5cm}|p{9cm}|}
\hline
\textbf{Функция} & \textbf{Назначение} \\
\hline
\texttt{pthread\_create} & Создание потока \\
\hline
\texttt{pthread\_join} & Ожидание завершения потока \\
\hline
\texttt{pthread\_exit} & Завершение потока \\
\hline
\texttt{pthread\_mutex\_lock} & Захват мьютекса \\
\hline
\texttt{pthread\_mutex\_unlock} & Освобождение мьютекса \\
\hline
\texttt{pthread\_cond\_wait} & Ожидание условной переменной \\
\hline
\texttt{pthread\_cond\_signal} & Сигнал условной переменной \\
\hline
\end{tabular}
\end{center}

\subsection{Минимальный пример POSIX Threads}

\begin{lstlisting}[language=C]
#include <stdio.h>
#include <pthread.h>

void* worker(void* arg) {
    int id = *(int*)arg;
    printf("Hello from thread %d\n", id);
    return NULL;
}

int main() {
    pthread_t t;
    int id = 1;

    pthread_create(&t, NULL, worker, &id);
    pthread_join(t, NULL);

    return 0;
}
\end{lstlisting}

Компиляция:

\begin{lstlisting}
gcc thread.c -pthread -o thread
\end{lstlisting}

\subsection{Мьютексы}

\textbf{Мьютекс} --- объект взаимного исключения, позволяющий только одному потоку входить в критическую секцию.

Пример:

\begin{lstlisting}[language=C]
#include <pthread.h>

int counter = 0;
pthread_mutex_t mutex = PTHREAD_MUTEX_INITIALIZER;

void* worker(void* arg) {
    for (int i = 0; i < 100000; i++) {
        pthread_mutex_lock(&mutex);
        counter++;
        pthread_mutex_unlock(&mutex);
    }
    return NULL;
}
\end{lstlisting}

\subsection{Условные переменные}

\textbf{Условная переменная} позволяет потоку ожидать наступления некоторого условия.

Типичный шаблон:

\begin{lstlisting}[language=C]
pthread_mutex_lock(&mutex);

while (!condition) {
    pthread_cond_wait(&cond, &mutex);
}

/* condition is true */

pthread_mutex_unlock(&mutex);
\end{lstlisting}

Важно использовать \texttt{while}, а не \texttt{if}, так как возможны ложные пробуждения.

\subsection{Семафоры POSIX}

Семафор --- счётчик, используемый для управления доступом к ограниченному числу ресурсов.

Основные функции:
\begin{itemize}
    \item \texttt{sem\_init};
    \item \texttt{sem\_wait};
    \item \texttt{sem\_post};
    \item \texttt{sem\_destroy}.
\end{itemize}

Пример:

\begin{lstlisting}[language=C]
#include <semaphore.h>

sem_t sem;

sem_init(&sem, 0, 1);

sem_wait(&sem);
/* critical section */
sem_post(&sem);

sem_destroy(&sem);
\end{lstlisting}

\subsection{Межпроцессное взаимодействие POSIX}

POSIX предоставляет несколько механизмов IPC:
\begin{itemize}
    \item каналы \texttt{pipe};
    \item именованные каналы FIFO;
    \item очереди сообщений;
    \item разделяемую память;
    \item семафоры;
    \item сигналы;
    \item сокеты.
\end{itemize}

\subsection{Пример канала POSIX}

\begin{lstlisting}[language=C]
#include <unistd.h>
#include <stdio.h>
#include <string.h>

int main() {
    int fd[2];
    char buffer[100];

    pipe(fd);

    if (fork() == 0) {
        close(fd[0]);
        write(fd[1], "hello", 6);
        close(fd[1]);
    } else {
        close(fd[1]);
        read(fd[0], buffer, sizeof(buffer));
        printf("Received: %s\n", buffer);
        close(fd[0]);
    }

    return 0;
}
\end{lstlisting}

\section{Сравнение PVM, MPI, OpenMP и POSIX}

\begin{center}
\begin{tabular}{|p{3cm}|p{3cm}|p{3cm}|p{3cm}|p{3cm}|}
\hline
\textbf{Критерий} & \textbf{PVM} & \textbf{MPI} & \textbf{OpenMP} & \textbf{POSIX} \\
\hline
Основная модель & Передача сообщений & Передача сообщений & Общая память & ОС-интерфейс \\
\hline
Тип системы & Распределённая, гетерогенная & Распределённая, HPC & Общая память & UNIX-подобные ОС \\
\hline
Единица выполнения & Задача PVM & MPI-процесс & Поток & Процесс/поток \\
\hline
Стандартизация & Слабее MPI & Высокая & Высокая & Высокая \\
\hline
Создание процессов & Динамическое через \texttt{pvm\_spawn} & Обычно через \texttt{mpirun}; есть динамические процессы & Потоки создаёт runtime & \texttt{fork}, \texttt{pthread\_create} \\
\hline
Обмен данными & Сообщения & Сообщения & Общие переменные & IPC, память, сокеты \\
\hline
Типичная область & Распределённые гетерогенные сети & Кластеры и суперкомпьютеры & Многоядерные узлы & Системное программирование \\
\hline
\end{tabular}
\end{center}

\section{Типовые проблемы управления процессами}

\subsection{Взаимная блокировка}

\textbf{Взаимная блокировка} возникает, когда процессы бесконечно ожидают друг друга.

Пример в MPI:

\begin{lstlisting}[language=C]
/* Process 0 */
MPI_Recv(&x, 1, MPI_INT, 1, 0, MPI_COMM_WORLD, &status);
MPI_Send(&y, 1, MPI_INT, 1, 0, MPI_COMM_WORLD);

/* Process 1 */
MPI_Recv(&y, 1, MPI_INT, 0, 0, MPI_COMM_WORLD, &status);
MPI_Send(&x, 1, MPI_INT, 0, 0, MPI_COMM_WORLD);
\end{lstlisting}

Оба процесса сначала ожидают получение сообщения, которое никто не отправляет.

Способы устранения:
\begin{itemize}
    \item изменить порядок операций;
    \item использовать \texttt{MPI\_Sendrecv};
    \item использовать неблокирующие операции;
    \item проектировать протокол обмена без циклического ожидания.
\end{itemize}

\subsection{Гонка данных}

Гонка данных характерна для многопоточных программ. Она возникает при некорректном доступе к общим данным.

Способы устранения:
\begin{itemize}
    \item мьютексы;
    \item атомарные операции;
    \item редукции;
    \item критические секции;
    \item минимизация общего состояния.
\end{itemize}

\subsection{Дисбаланс нагрузки}

\textbf{Дисбаланс нагрузки} возникает, когда одни процессы завершают работу раньше других и простаивают.

Способы борьбы:
\begin{itemize}
    \item динамическое распределение задач;
    \item разбиение на мелкие блоки;
    \item использование планирования \texttt{dynamic} в OpenMP;
    \item master-worker схема в MPI или PVM;
    \item профилирование программы.
\end{itemize}

\section{Алгоритм master-worker}

\subsection{Описание}

Алгоритм master-worker применяется для динамического распределения независимых задач между рабочими процессами.

\textbf{Master} хранит очередь задач и выдаёт их свободным рабочим процессам.

\textbf{Worker} получает задачу, выполняет её и возвращает результат.

\subsection{Шаги алгоритма}

\begin{enumerate}
    \item Главный процесс формирует список задач.
    \item Рабочие процессы отправляют запросы на получение задачи.
    \item Главный процесс выдаёт задачу свободному рабочему.
    \item Рабочий процесс выполняет задачу.
    \item Рабочий процесс отправляет результат главному процессу.
    \item Если задачи остались, рабочему выдаётся новая задача.
    \item Если задач нет, главный процесс отправляет сигнал завершения.
    \item Все процессы завершаются.
\end{enumerate}

\subsection{Иллюстрация}

\[
\text{Очередь задач: } Q = \{q_1,q_2,\ldots,q_m\}
\]

\[
\text{Master} \rightarrow \text{Worker}_i : q_k
\]

\[
\text{Worker}_i \rightarrow \text{Master} : r_k
\]

\subsection{Минимальная MPI-схема}

\begin{lstlisting}[language=C]
if (rank == 0) {
    for (int task = 1; task < size; task++) {
        MPI_Send(&task, 1, MPI_INT, task, 0, MPI_COMM_WORLD);
    }
} else {
    int task;
    MPI_Recv(&task, 1, MPI_INT, 0, 0, MPI_COMM_WORLD,
             MPI_STATUS_IGNORE);
    /* выполнить task */
}
\end{lstlisting}

\section{Гибридное программирование MPI + OpenMP}

\subsection{Идея}

В современных кластерах каждый узел содержит несколько ядер. Поэтому часто используют:
\begin{itemize}
    \item MPI для обмена между узлами;
    \item OpenMP для распараллеливания внутри узла.
\end{itemize}

Схема:

\[
\text{Узел}_1:
\begin{cases}
\text{MPI-процесс}_0 \rightarrow T_0,T_1,\ldots,T_k
\end{cases}
\]

\[
\text{Узел}_2:
\begin{cases}
\text{MPI-процесс}_1 \rightarrow T_0,T_1,\ldots,T_k
\end{cases}
\]

\subsection{Минимальный пример}

\begin{lstlisting}[language=C]
#include <mpi.h>
#include <omp.h>
#include <stdio.h>

int main(int argc, char **argv) {
    int rank, size;

    MPI_Init(&argc, &argv);

    MPI_Comm_rank(MPI_COMM_WORLD, &rank);
    MPI_Comm_size(MPI_COMM_WORLD, &size);

#pragma omp parallel
    {
        int tid = omp_get_thread_num();
        printf("MPI rank %d/%d, OpenMP thread %d\n",
               rank, size, tid);
    }

    MPI_Finalize();
    return 0;
}
\end{lstlisting}

\section{Критерии выбора средства}

\subsection{Когда использовать MPI}

MPI целесообразен, если:
\begin{itemize}
    \item программа должна работать на кластере;
    \item память распределена между узлами;
    \item требуется масштабирование на большое число процессов;
    \item нужен контроль над обменами;
    \item важна высокая производительность.
\end{itemize}

\subsection{Когда использовать OpenMP}

OpenMP целесообразен, если:
\begin{itemize}
    \item программа выполняется на одном многоядерном узле;
    \item данные помещаются в общую память;
    \item нужно быстро распараллелить циклы;
    \item требуется сохранить структуру последовательной программы.
\end{itemize}

\subsection{Когда использовать POSIX}

POSIX используется, если:
\begin{itemize}
    \item требуется системное программирование;
    \item нужно явно управлять процессами или потоками;
    \item требуется переносимость между UNIX-подобными системами;
    \item нужны низкоуровневые IPC-механизмы.
\end{itemize}

\subsection{Когда использовать PVM}

PVM может быть полезен, если:
\begin{itemize}
    \item требуется объединить гетерогенные машины в виртуальную вычислительную систему;
    \item нужна динамическая модель добавления и удаления узлов;
    \item изучаются исторические и концептуальные основы распределённых вычислений.
\end{itemize}

\section{Заключение}

Средства управления процессами в параллельных и распределённых системах различаются по уровню абстракции и целевой архитектуре.

\begin{itemize}
    \item \textbf{PVM} предоставляет модель виртуальной параллельной машины и удобен для гетерогенных распределённых систем.
    \item \textbf{MPI} является основным стандартом высокопроизводительного программирования с передачей сообщений.
    \item \textbf{OpenMP} предназначен для параллельного программирования в системах с общей памятью и удобен для распараллеливания циклов.
    \item \textbf{POSIX} задаёт переносимые системные интерфейсы управления процессами, потоками и межпроцессным взаимодействием.
\end{itemize}

На практике эти средства часто комбинируются. Например, MPI используется для взаимодействия между узлами кластера, OpenMP --- для распараллеливания внутри узла, а POSIX --- для низкоуровневого управления процессами, потоками, файлами и сигналами.

\end{document}
